\subsection{区间的长度}

定义 3. 以 \(a\) 与 \(b\) （ \(a \leqslant b\) ）为端点的有界区间，其长度定义为 \(b - a\) 。

因此每个包含多于一个点的有界区间长度 \(> 0\) 。若 \(a \leqslant b\) ，则四个区间 \([a, b]\) 、\(]a, b[\) 、\([a, b[\) 与 \(]a, b[\) 都有相同的长度。以 \(a + c\) 与 \(b + c\) 为端点的区间与以 \(a\) 与 \(b\) 为端点的区间有相同的长度；换句话说，区间的长度在平移下不变。

若 \(a \leqslant c \leqslant d \leqslant b\) ，我们有 \(d - c \leqslant b - a\) 。因此若有界区间 \(I\) 含于有界区间 \(I'\) ，则 \(I\) 的长度小于或等于 \(I'\) 的长度。

若 \(n\) 个两两不相交的开区间 \(I_1, I_2, \ldots, I_n\) 含于区间 \([a, b] (a < b)\) ，则对 \(n\) 用归纳法容易看出：若 \(I_k = ]c_k, d_k[\) ，则存在指标的一个置换 \(\sigma\) （ \(k (1 \leqslant k \leqslant n)\) ），使得 \(d_{\sigma(k)} \leqslant c_{\sigma(k+1)}\) 对 \(1 \leqslant k \leqslant n - 1\) 成立。由此立即推出，诸区间 \(I_k\) 的长度之和至多等于 \([a, b]\) 的长度，且等号成立当且仅当 \(c_{\sigma(1)} = a, d_{\sigma(n)} = b\) 且 \(d_{\sigma(k)} = c_{\sigma(k+1)}\) 对 \(1 \leqslant k \leqslant n - 1\) 成立。

